On dit qu'une droite du plan est \emph{ensoleillée} lorsqu'elle n'est
parallèle ni à l'axe des abscisses, ni à l'axe des ordonnées, ni à la
droite d'équation $x + y = 0$.

Soit $n \geqslant 3$ un entier donné. Trouver tous les entiers
$k \geqslant 0$ pour lesquels il existe $n$ droites du plan, deux à
deux distinctes, telles que
\begin{itemize}
  \item pour tous les entiers $a \geqslant 1$ et $b \geqslant 1$ tels
  que $a + b \leqslant n + 1$, le point de coordonnées $(a, b)$
  appartient à au moins une de ces $n$ droites, et
  \item exactement $k$ de ces $n$ droites sont ensoleillées.
\end{itemize}
